% -----------------------------------------------------------------------------
% Supplied PDMS mixing tutorial
% Source: PDMS_Mixing_Tutorial.docx
% Wording and source-image order are preserved.
% -----------------------------------------------------------------------------

\clearpage
\hypertarget{q4-pdms-mixing-tutorial}{}
\sectiontitle{Q4. PDMS Mixing Tutorial (Polydimethylsiloxane)}
\qtwosubsection{Practical mixing procedure}

\begingroup
\fontsize{9.0}{10.45}\selectfont
\newcommand{\pdmshead}[1]{\par\vspace{0.48em}{\bfseries #1}\par\vspace{0.15em}}

This tutorial describes the general process for mixing PDMS silicone elastomer
for molding, coatings, and laboratory applications.

\pdmshead{Materials}
\begin{itemize}[leftmargin=1.45em,itemsep=0.15em,topsep=0.25em]
  \item PDMS base (silicone prepolymer)
  \item Curing agent (crosslinker)
  \item Mixing cup and spatula
  \item Balance or measuring tools
  \item Vacuum desiccator (for degassing)
  \item Oven or hot plate for curing
\end{itemize}

\pdmshead{Step 1: Measure the Components}
We used Sylgard 184 mixing ratio which is 10:1 by weight (base:curing agent).
Example: 10 g PDMS base mixed with 1 g curing agent. Different ratios can change
the softness, curing time, and mechanical properties.

\pdmshead{Step 2: Mix Thoroughly}
Add the curing agent to the PDMS base and mix slowly for 5--10 minutes. Scrape
the sides and bottom of the container during mixing. Avoid vigorous stirring
because it introduces air bubbles.

\pdmshead{Step 3: Degas the PDMS}
Place the mixed PDMS in a vacuum chamber. Apply vacuum until trapped air bubbles
expand and disappear. This usually takes 15--60 minutes depending on the volume
and viscosity. Release the vacuum slowly to prevent overflow.

\pdmshead{Step 4: Pour or Coat}
Transfer the degassed PDMS onto the mold or substrate carefully. Pour slowly to
reduce the chance of trapping new air bubbles.

\pdmshead{Step 5: Cure the PDMS}
Typical Sylgard 184 curing conditions are 60--80~$^{\circ}$C for 1--2 hours.
PDMS can also cure at room temperature, but this may require about 24 hours or
longer depending on conditions.

\pdmshead{Step 6: Demold and Finish}
Allow the PDMS to cool after curing. Carefully remove it from the mold and trim
edges or excess material as required.

\pdmshead{Important Tips}
\begin{itemize}[leftmargin=1.45em,itemsep=0.15em,topsep=0.25em]
  \item More curing agent generally produces a harder and faster-curing PDMS.
  \item Less curing agent generally produces softer and more flexible PDMS.
  \item Proper degassing is important for removing bubbles.
  \item Use clean tools because contamination may affect curing.
\end{itemize}
\endgroup

\clearpage
\sectiontitle{Q4. PDMS Mixing Tutorial -- Laboratory Images}
\begin{center}
  \centering
  \includegraphics[angle=270,origin=c,width=0.47\linewidth,height=0.36\textheight,keepaspectratio]{assets/pdms_mixing/01_balance_setup.jpeg}\hfill
  \includegraphics[angle=270,origin=c,width=0.47\linewidth,height=0.36\textheight,keepaspectratio]{assets/pdms_mixing/02_component_handling.jpeg}

  \vspace{0.8em}
  \includegraphics[angle=270,origin=c,width=0.47\linewidth,height=0.36\textheight,keepaspectratio]{assets/pdms_mixing/03_curing_agent_mass.jpeg}\hfill
  \includegraphics[angle=270,origin=c,width=0.47\linewidth,height=0.36\textheight,keepaspectratio]{assets/pdms_mixing/04_pdms_base_mass.jpeg}
\end{center}

\clearpage
\sectiontitle{Q4. PDMS Mixing Tutorial -- Laboratory Images}
\begin{center}
  \centering
  \includegraphics[angle=270,origin=c,width=0.47\linewidth,height=0.70\textheight,keepaspectratio]{assets/pdms_mixing/05_manual_mixing.jpeg}\hfill
  \includegraphics[angle=270,origin=c,width=0.47\linewidth,height=0.70\textheight,keepaspectratio]{assets/pdms_mixing/06_vacuum_oven.jpeg}
\end{center}
